% exp11 (from-scratch matched-budget LM training) discussion blocks.
% pick ONE of the three blocks below and toggle its \iftrue / \iffalse.
%
% decision rule: read SUMMARY.md and pull final val_loss for swiglu (S),
% default-init tucker (T_default), and diagonal-bias-init tucker (T_hill).
%
%   T_default = 5.116  (results/exp11/tucker_seed0/loss_log.json)
%   T_hill_v1 = 4.772  (results/exp11_hc/tucker_seed0/loss_log.json)
%   T_hill_v2 = 4.770  (results/exp11_hc_v2/tucker_seed0/loss_log.json,
%                       v1 + --tucker_core_lr_scale 2.0)
%   S         = 4.758  (results/exp11/swiglu_seed0/loss_log.json)
%
%   delta_hill_v1_swiglu = T_hill_v1 - S = +0.014  (~1.5% perplexity gap)
%   delta_hill_v2_swiglu = T_hill_v2 - S = +0.012  (~1.3% perplexity gap)
%   delta_hill_v2_default = T_default - T_hill_v2 = +0.346  (~30% perplexity
%                                                            reduction from
%                                                            the hill-climb)
%
% framing rule (per CLAUDE_CODE_PROMPT §Gap 3):
%   (a) hill-climb wins        T_hill < S - 0.005
%   (b) hill-climb close/lose  T_hill >= S, T_default > T_hill   <-- current
%   (c) cut exp11              both tucker variants lose decisively
%                              (T_hill > S + 0.05)
%
% with current numbers framing (b) applies: hill-climb closed essentially the
% entire optimization gap (diagonal_bias_init halves the perplexity gap
% relative to swiglu) but does not strictly beat it. if the v2 hill-climb
% (--tucker_core_lr_scale 2.0, results/exp11_hc_v2/) lands below 4.753, switch
% to framing (a) and update the numbers; if v2 also lands close-but-behind,
% keep framing (b) and append the v2 number to the table.

% ── Framing (a): hill-climb wins ─────────────────────────────────────────
\iffalse
At matched parameter count ($d{=}512$, $L{=}8$, $\sim$50M params,
$r{=}s{=}128$ for Tucker matched against $m{=}1493$ for SwiGLU), Tucker
matches or exceeds SwiGLU's validation loss at 100M FineWeb-Edu tokens
when $C$ is initialized with a superdiagonal bias so that the layer starts
in the SwiGLU-equivalent subspace and is free to deviate into off-diagonal
interactions only when training prefers it
(\(T_{\textsc{hill}}{=}T_HILL\) vs.\ \(S{=}S_VAL\); Figure~\ref{fig:lm_loss_curves}).
Random Tucker initialization does not reach this regime at our scale
(\(T_{\textsc{default}}{=}T_DEFAULT\)), suggesting the representational gap
of Theorem~\ref{thm:separation} is real but only exploitable end-to-end
under a SwiGLU-aware warm start. Stable-rank analysis of the resulting
$V_j$ confirms the hill-climbed model uses cross-channel interactions of
mean rank $\approx 17$ per gate (Section~\ref{sec:exp12}), well above the
rank-one-per-gate ceiling that an aligned SwiGLU of width $r$ would impose.
\fi

% ── Framing (b): hill-climb close, distillation framing ─────────────────
\iftrue
At matched parameter count ($d{=}512$, $L{=}8$, $\sim$50M params,
$r{=}s{=}128$ for Tucker matched against $m{=}1493$ for SwiGLU), SwiGLU
slightly outperforms Tucker on 100M FineWeb-Edu tokens
(val\_loss $S{=}4.758$ vs.\ $T_{\textsc{hill}}{=}4.772$, a $\sim$$1.5\%$
perplexity gap; Figure~\ref{fig:lm_loss_curves}). Random Tucker
initialization is much worse ($T_{\textsc{default}}{=}5.116$), and biasing
$C$ toward its superdiagonal at initialization closes nearly all of that
gap, indicating that the optimization landscape at our scale favors the
diagonal-restricted parameterization but the residual gap is small once
the optimizer is given a SwiGLU-equivalent starting point. The
representational gap that Theorem~\ref{thm:separation} predicts is
independently confirmed by distillation (Section~\ref{sec:exp14b}): when
the teacher is a trained Tucker layer, a Tucker student matches it
roughly $22\%$ more accurately than a parameter-matched SwiGLU student
at the same $r{=}s{=}128$ matched-budget configuration. Stable-rank
analysis of the hill-climbed Tucker (Section~\ref{sec:exp12}) finds mean
$\mathrm{rank}(V_j) \approx 17$ per gate, well above the rank-one-per-gate
ceiling an aligned SwiGLU of width $r$ would impose, so the model is
demonstrably exploiting cross-channel structure that Theorem~1 forbids
SwiGLU from representing. The diagonal restriction therefore does cost
representational capacity at the layer level (Theorem~\ref{thm:separation}
is empirically binding), but the matching optimization gap absorbs most of
the resulting end-to-end advantage at workshop training scale; whether
the representational advantage materializes end-to-end at frontier scale
is open.
\fi

% ── Framing (c): cut exp11 ───────────────────────────────────────────────
\iffalse
We do not include a from-scratch language-modeling comparison in the main
text: at our matched-budget training scale ($\sim$50M params, 100M tokens)
the optimization gap between random-init Tucker and SwiGLU dominates the
representational signal that Theorem~\ref{thm:separation} surfaces. The
diagonal-projection cost (Section~\ref{sec:exp13}) and the matched-budget
distillation comparison (Section~\ref{sec:exp14b}) measure the
representational signal directly without confounding it with optimization,
and we leave a careful study of Tucker training dynamics --- including
whether the diagonal-bias initialization gap closes with longer training,
larger models, or different optimizers --- to future work.
\fi
